%%%%%%%%%%%%%%%%%%%%%%%%%%%%%%%%%%%%%%%%%%%%%%%%%%%%%%%%%%%%%%%%%%%%%%
% Câu chưa có mã dạng — cần gắn %[..] rồi chuyển vào đúng file dạng
% Số câu: 1  (sinh tự động bởi tach_cau.py; từ nay thầy sửa trực tiếp file này)
%%%%%%%%%%%%%%%%%%%%%%%%%%%%%%%%%%%%%%%%%%%%%%%%%%%%%%%%%%%%%%%%%%%%%%

%%%%%%%%%%%%%%%% ĐÚNG/SAI %%%%%%%%%%%%%%%%

% nguồn: Ok_dungsai_lop11_Lan2
\begin{ex}%[ma:37fe55a7]
	~%[1H4N6-1] 
	 \choiceTF
		{\True Hình biểu diễn của một hình bình hành là một  hình bình hành hoặc là một đoạn thẳng}
		{Hình biểu diễn của một hình chữ nhật là một hình chữ  nhật hoặc là một đoạn thẳng}
		{Hình  biểu diễn của một hình vuông là một hình vuông hoặc là một đoạn thẳng}
		{Hình biểu diễn của một hình thoi là một hình thoi hoặc là một đoạn  thẳng}
	 \loigiai{
		\begin{itemchoice}
			\itemch Hình biểu diễn của  một hình bình hành là một hình bình hành hoặc là một đoạn thẳng.
			\itemch Điều chỉnh lại cho \lq\lq Hình biểu diễn của  một hình chữ nhật là một hình bình hành hoặc là một đoạn thẳng\rq\rq.
			\itemch Điều chỉnh  lại cho \lq\lq Hình biểu diễn của một hình vuông là một hình bình hành hoặc là một đoạn thẳng\rq\rq.
			\itemch Điều chỉnh lại cho \lq\lq Hình biểu  diễn của một hình thoi là một hình bình hành hoặc là một đoạn thẳng\rq\rq.
		\end{itemchoice}
	} 
\end{ex}
